\cotefile{29}
\scanpage{25}
% archive sheet number: 24
obvious when $P=G'_S$ with the right regular action of $G'$, and the left action of $G$ is given by a monomorphism $G\quad G'$; one reduces to this particular case by étale localization.

One may propose to normalize the notion of ramification datum, to eliminate the arbitrariness in the choice of an $\underline{R}$. But then \unsure{\deleted{o}it|\deleted{i}ol|\deleted{b}it|\deleted{d}it|\deleted{i}bl|\deleted{i}dl} seems that one is infallibly led to take as category $\underline{C}$ the category $\underline{\mathrm{Rev}}^{\underline{R}}$ itself which one proposes \unsure{t\deleted{'}o|t\deleted{e}o} construct and to study ! (Compare with the analogous situation for the relations between sites and topoi, the auxiliary object being the site, the intrinsic object of geometric naturenbeing the topos associated with this site).

2. \uline{Properties of the category} $\operatorname{Rev}^{\underline{R}}(S)$. \uline{The fundamental group} $ {}^{\underline{R}}_1(S,\quad)$.

One observes immediately that the category $\operatorname{Rev}^{\underline{R}}(S)$ inherits all the properties of ex\unsure{a|\deleted{e}a|\deleted{a}e}\unsure{s|c|\deleted{s}c|\deleted{c}s}tness \deleted{already} established for the category of étalecoverings of a given schemes, and some of which are stated in Exp V number 4 (conditions G1 to G3). To be precise, one can say that in the category $\operatorname{Rev}^{\underline{R}}(S)$, finite projective limits and finite inductive limits exist, and satisfy the same properties of exactness as in the category of sets (such as the exactness of the base change functor in this category); \deleted{moreover\unsure{;|,}} \inserted{in particular} every morphism of the category factors as the product of an epimorphism and a mnomorphism, \inserted{the epimorphisms and monomorphisms are \deleted{\unsure{s}tri} effective; finally} every monomorphism is an isomorphism with a direct summand. One will take note that the functor ``geometric realization'' on $\operatorname{Rev}^{\underline{R}}(S)$ does commute with finite inductive limits (and \deleted{consequently with\unsure{x}} in particular with finite sums), but not in general with finite projective limits, and in particular, not with products of two objects.

To apply the general Galois theory of V 4, one must dis-
